\VerdierPhysicalPageAnchor{0035}%
\VerdierSemanticAnchor{ch1-1-1-1-fiber-category-cont}%
donc la sous-catégorie de $\mathcal C$ ayant même ensemble d'objets que
$\mathcal C$ et dont les morphismes sont les morphismes de $\mathcal C$ de
degré $e$.

\VerdierSemanticAnchor{ch1-1-1-1-graded-functor}
Soient $(\mathcal C,G)$ et $(\mathcal C',G)$ deux catégories graduées de
type $G$. On appelle \emph{foncteur gradué} de $(\mathcal C,G)$ dans
$(\mathcal C',G)$ un foncteur cartésien
$F$ au-dessus de $\mathcal G$. On désigne par
$F_e\colon\mathcal C_e\to\mathcal C'_e$ le foncteur obtenu par restriction
aux catégories fibres.
\VerdierSemanticAnchor{ch1-1-1-1-compatible-morphism}
Enfin, un morphisme de foncteurs gradués $m\colon F\to F'$ est dit
\emph{compatible avec la graduation} si c'est un
$\mathcal G$-morphisme de foncteurs.

\VerdierSemanticAnchor{ch1-1-1-2}
\noindent\textbf{1.1.2.}\quad
Soient $\mathcal C$ une catégorie graduée de type $G$ et $g$ un élément de
$G$. Pour tout couple $X,Y$ d'objets de $\mathcal C$, on désigne par
\VerdierSemanticAnchor{ch1-1-1-2-homogeneous-hom}%
$\operatorname{Hom}^{g}(X,Y)$ l'ensemble des morphismes de $X$ dans $Y$ de
degré $g$. On a alors :
\VerdierSemanticAnchor{ch1-1-1-2-eq1}
\begingroup
\makeatletter
\tagsleft@true
\let\veqno\@@leqno
\makeatother
\begin{equation}
  \operatorname{Hom}(X,Y)
  =\coprod_{g\in G}\operatorname{Hom}^{g}(X,Y).
  \tag{1.1.2.1}
\end{equation}
\endgroup

\VerdierSemanticAnchor{ch1-1-1-2-graded-hom-set}
L'ensemble $\operatorname{Hom}(X,Y)$ est donc muni d'une structure
d'ensemble gradué de type $G$. Le degré étant fonctoriel, on a pour tout
morphisme $m\colon X\to Y$ et tout morphisme $n\colon Y\to Z$ :
\VerdierSemanticAnchor{ch1-1-1-2-eq2}
\begin{equation}
  \deg(n)\deg(m)=\deg(nm).
  \tag{1.1.2.2}
\end{equation}

\VerdierSemanticAnchor{ch1-1-1-2-fibered-criterion}
Les formules précédentes ne font que traduire le fait que $\mathcal C$ est
une catégorie au-dessus de $\mathcal G$. Exprimons maintenant que la
catégorie $\mathcal C$ est une catégorie fibrée au-dessus de $\mathcal G$.
Il suffit pour cela d'exprimer que tout objet de $\mathcal C$ est but d'un
morphisme hypercartésien de degré donné [19]. Or on vérifie immédiatement
que les morphismes hypercartésiens de $\mathcal C$ sont les isomorphismes
de $\mathcal C$, et par suite que la catégorie $\mathcal C$ est fibrée si
et seulement si elle possède la propriété suivante :

\VerdierSemanticAnchor{ch1-1-1-2-eq3}
\noindent\textnormal{(1.1.2.3)}\quad
\begin{minipage}[t]{0.78\textwidth}
\emph{Pour tout élément $g$ de $G$ et tout objet $X$ de $\mathcal C$, il
existe un isomorphisme de but $X$ et de degré $g$.}
\end{minipage}
\par\medskip

\VerdierSemanticAnchor{ch1-1-1-2-degree-preserving-functors}
Soient maintenant $(\mathcal C,G)$ et $(\mathcal C',G)$ deux catégories
graduées de type $G$ et $F\colon\mathcal C\to\mathcal C'$ un foncteur
gradué. Le foncteur $F$ est d'abord un $\mathcal G$-foncteur et, par suite,
il conserve le degré des morphismes. Il transforme de plus les morphismes
cartésiens de $\mathcal C$ en morphismes cartésiens de $\mathcal C'$, ce
qui est automatique car les morphismes cartésiens de $\mathcal C$ sont des
isomorphismes. Les foncteurs gradués de $\mathcal C$ dans $\mathcal C'$
sont donc simplement \emph{les foncteurs de $\mathcal C$ dans
$\mathcal C'$ qui conservent le degré}.

\VerdierSemanticAnchor{ch1-1-1-2-degree-e-morphism}
Enfin, un morphisme de foncteurs gradués $m\colon F\to F'$ compatible avec
la graduation est simplement un morphisme de foncteurs de degré $e$,
\emph{i.e.} la donnée pour tout objet $X$ de $\mathcal C$ d'un morphisme :
\VerdierSemanticAnchor{ch1-1-1-2-component}
\[
  m(X)\colon F(X)\longrightarrow F'(X)
\]
